\documentclass[10pt]{article}
\usepackage[margin=20mm]{geometry}
\usepackage{amsmath,amssymb,mathtools,fontspec,unicode-math,xurl,hyperref}
\IfFontExistsTF{DejaVu Serif}{\setmainfont{DejaVu Serif}}{\setmainfont{Latin Modern Roman}}
\setmathfont{latinmodern-math.otf}
\setlength{\parindent}{0pt}\setlength{\parskip}{6pt}\sloppy
\title{Technical Verification of
Unramified Grothendieck–Serre for Isotropic Groups}
\author{AI-assisted technical review}
\date{}
\begin{document}
\maketitle
Abstract: The finite-field reembedding criterion is examined together with its uses in the proof of unramified Grothendieck–Serre. Under the standard meaning of rational points, Definition 2.5 does not give a sufficient condition for Lemma 2.6. An explicit connected scheme over \(\mathbb F_2\) satisfies every stated finite-extension point-count inequality but cannot embed in the prescribed open affine line. The corrected finite-field criterion counts closed points of each exact residue degree. Two special applications in the manuscript satisfy this stronger condition; the external cover used in the main proof has not been independently checked here. The finding concerns the stated reembedding justification, not a counterexample to the main torsor theorem. Scientific adjudication remains pending.\par
\section*{Paper identity and scope}
The reviewed source is Kęstutis Česnavičius and Roman Fedorov, Unramified Grothendieck–Serre for isotropic groups, arXiv:2311.08660v3, 21 November 2025 [CF]. The PDF carries an internal date of 24 November 2025. The acquired identity record associates it with Journal of the European Mathematical Society, DOI https://doi.org/10.4171/jems/1725, and describes it as the final author manuscript with journal numbering updated. No publisher byte equivalence is asserted.\par
All 24 pages of the inherited text extraction were read, including the bibliography. This does not constitute verification of all proofs or cited works. The substantive checks concern the finite-field criterion, its concrete applications, and selected reductions in Theorem 3.5 and the concluding proofs. Actual images of pages 1, 2, 4, 8, 9, 11, 15, 16, 18, 19, and 20 were inspected. PDF and printed page numbers agree.\par
\section*{Hypotheses and distinctions between the main claims}
Theorem 1.2, restated as Theorem 4.3, assumes a Noetherian semilocal ring \(R\), flat and geometrically regular over a Dedekind ring \(\mathcal O\), and a reductive \(R\)-group \(G\) whose adjoint quotient is totally isotropic. It asserts \[\ker\bigl(H^1(R,G)\longrightarrow H^1(\mathop{\mathrm{Frac}}(R),G)\bigr)=\{*\}\] [CF, p. 2]. Section 1.3 defines total isotropy factorwise through the finite étale and Weil restriction decomposition of the adjoint group: the simple factors admit proper parabolic subgroups Zariski locally, containing no whole fiber [CF, p. 4]. No assertion for arbitrary ramified \(R\) or arbitrary anisotropic factors is substituted for this statement.\par
Theorem 3.5 has a different scope: a torsor under any reductive group over a semilocal \(A\) is sectionwise constant on \(\mathbb P_A^d\). The stronger affine-line triviality in Theorem 4.2 retains total isotropy, although it allows arbitrary \(A\). Proposition 5.1 lifts generically trivial torsors to the simply connected derived group under the unramified ring hypotheses. The issue below is confined to an auxiliary reembedding criterion used in these arguments.\par
\section*{A finite-field counterexample to the printed criterion}
Problem 1 (Rational-point inequalities do not ensure closed embeddings).  With \(k'\)-points understood as \(\operatorname{Hom}_k(\mathop{\mathrm{Spec}}k',-)\), the condition in Definition 2.5 is insufficient for the sufficiency direction of Lemma 2.6.\par
\subsection*{Original claim}
Definition 2.5 calls the finite-field obstruction absent if, over every finite residue field \(k\) of the base and every finite extension \(k'/k\), the number of \(k'\)-points of the source does not exceed that of the target. Lemma 2.6 states that a finite \(k\)-scheme \(Y\) embeds as a closed subscheme of a nonempty open \(V\subset\mathbb A_k^1\) if and only if this condition holds and \(Y\) is a closed subscheme of some smooth \(k\)-curve [CF, p. 8]. The first sentence of its proof on page 9 uses this condition to choose an embedding of the closed point of a connected \(Y\). The actual page images confirm this wording.\par
\subsection*{Parameters and verification of all hypotheses}
Set \[k=\mathbb F_2,\qquad P(t)=t^2+t+1,\qquad
 Y=\mathop{\mathrm{Spec}}\mathbb F_4,\qquad
 V=\mathop{\mathrm{Spec}}\mathbb F_2[t,1/P(t)].\] The polynomial \(P\) has no root in \(\mathbb F_2\), so it is irreducible of degree two. Thus \(Y\) is the closed subscheme of the smooth connected curve \(C=\mathbb A_{\mathbb F_2}^1\) cut out by \(P\). It is finite and connected. The target \(V=D(P)\) is a nonempty open of the affine line; both 0 and 1 belong to it. The field \(k\) is itself a semilocal base with its sole residue field equal to \(k\).\par
Equation reference: eq:Ycounts, eq:Vcounts. For every \(m\ge1\), the finite extension of degree \(m\) is \(\mathbb F_{2^m}\), and \[\begin{aligned}
 \#Y(\mathbb F_{2^m})&=
 \begin{cases}
 2,&2\mid m,\\
 0,&2\nmid m,
 \end{cases}
 \\
 \#V(\mathbb F_{2^m})&=
 \begin{cases}
 2^m-2,&2\mid m,\\
 2^m,&2\nmid m.
 \end{cases}
 \end{aligned}\] Indeed, \(\mathbb F_4\) embeds in \(\mathbb F_{2^m}\) precisely when \(2\mid m\), in which case there are two embeddings. The same two conjugate roots of \(P\) are exactly the points removed from \(\mathbb A^1(\mathbb F_{2^m})\). If \(m\) is even, then \(m\ge2\) and \(2\le2^m-2\). If \(m\) is odd, the source has no points. Consequently every inequality in the printed definition holds. This proves the quantified condition for all finite extensions, rather than testing finitely many of them.\par
\subsection*{Failure of the asserted embedding}
Equation reference: eq:map. A closed embedding \(Y\hookrightarrow V\) would induce a surjective map \[
 \mathbb F_2[t,1/P(t)]\longrightarrow\mathbb F_4.\] Write \(a\) for the image of \(t\). If \(a\in\mathbb F_2\), the image of [eq:map] (\#eq:map) lies in \(\mathbb F_2\), so the map is not surjective. If \(a\notin\mathbb F_2\), then \(P(a)=0\), contradicting the fact that \(P(t)\) is inverted in the source ring. These cases exhaust \(\mathbb F_4\). Thus no closed embedding exists. Equivalently, the unique degree-two closed point of \(\mathbb A_{\mathbb F_2}^1\) has been deleted from \(V\).\par
This is an admissible counterexample to the printed sufficiency criterion. It does not construct a nontrivial generically trivial torsor under the hypotheses of the principal theorem.\par
\subsection*{Correct finite-field condition}
Equation reference: eq:degrees. Let \(c_d(T)\) denote the number of closed points of \(T\) with residue degree exactly \(d\) over \(\mathbb F_q\). A necessary condition for a closed embedding is \[
 c_d(Y)\le c_d(V)\qquad(d\ge1).\] The rational-point counts satisfy \[\#T(\mathbb F_{q^m})=\sum_{d\mid m}d\,c_d(T).\] Inequalities between these sums do not imply inequalities between their individual terms. In the example, \(c_2(Y)=1\) and \(c_2(V)=0\); the degree-one points of \(V\) mask that shortage in [eq:Vcounts] (\#eq:Vcounts).\par
For finite \(k\), condition [eq:degrees] (\#eq:degrees), together with the assumed smooth-curve embedding of \(Y\), is also sufficient. Each local component of \(Y\) has the form \(k_d[u]/(u^e)\), because a finite field is perfect and the curve is smooth. Choose distinct target points of \(V\) of the corresponding degrees. Their length-\(e\) infinitesimal neighborhoods have the same local form, so the components embed. A prescribed embedding on a closed subscheme \(Y_0\) can be retained: keep its target points and extend each truncated coordinate isomorphism. The nonzero linear coefficient of a uniformizer change remains nonzero at every larger truncation order.\par
This verifies the finite-field correction, including the extension aspect in that case. It is not a new proof of the infinite-field inseparable argument or of every subsequent application. If the intended meaning of \(k'\)-points was closed points having residue field \(k'\), the intended condition is [eq:degrees] (\#eq:degrees); that interpretation requires an explicit change from the usual rational-point meaning analyzed above.\par
\section*{Dependency, surviving applications, and limits}
The criterion is a substantive reembedding input. The proof of Theorem 4.3 first invokes a finite étale cover from the externally cited Lemma 6.1 of the paper’s reference [Čes22a], states the absence of a finite-field obstruction, and then applies Lemmas 2.6 and 2.7 to pass to an affine-line excision square [CF, p. 18]. The general criterion used in that inference is false as written. The affected reembedding justification is closely tied to the core proof, so the issue is assessed as Major under the claim-level severity convention.\par
This assessment has two qualifications. First, the external cover on page 18 has not been independently inspected here. It may satisfy the stronger exact-degree condition. The present finding does not assert that it fails, and does not refute Theorem 4.3. Second, two concrete uses in the manuscript already satisfy the stronger condition.\par
In Proposition 2.8 the fibers have been reduced to support on one rational point, and an additional disjoint rational point is introduced [CF, p. 11]. The affine line has at least two rational points over every field, including \(\mathbb F_2\). The required degree-one slots therefore exist, and no higher-degree slots are needed. This special application survives.\par
In the proof of Theorem 3.5, the construction \(Y=Y_N\sqcup Y_{N+1}\) on page 16 also meets [eq:degrees] (\#eq:degrees) for sufficiently large \(N\). Over a finite residue field \(\mathbb F_q\), write the finitely many components of \(Y_0\) as residue fields of degrees \(d_i\). The constructed covers have degrees \(d_iN\) and \(d_i(N+1)\), and a bounded number of components. If \(I_q(m)\) is the number of degree-\(m\) closed points of \(\mathbb G_m\), then for \(m>1\) \[mI_q(m)\ge q^m-\sum_{j=1}^{\lfloor m/2\rfloor}q^j,\] since every element of \(\mathbb F_{q^m}\) of degree less than \(m\) lies in a proper subfield of degree at most \(m/2\). Hence \(I_q(m)\to\infty\). The available exact-degree points eventually exceed the fixed component count. Semilocality gives only finitely many residue fields, so a single large \(N\) works for all of them. This verifies the relevant correction for the page-16 finite-field construction.\par
What remains for a complete correction of the affected main proof is to state [eq:degrees] (\#eq:degrees) and verify it for the particular cover and shrinking used on page 18. No author correction or adoption is established by the material consulted.\par
\section*{Other reductions checked}
The projective-line reduction in Theorem 3.5 respects the semilocal hypothesis [CF, p. 15]. Given two sections, choose in each closed fiber a third point different from both. Even \(\mathbb P^1(\mathbb F_2)\) has three points. Chinese remainders lift the finitely many choices to a unimodular vector over \(A\), hence an \(A\)-section. Two disjoint sections span a rank-two direct summand: their rank-two minors generate the unit ideal. Their span therefore gives a projective line. This checks the geometric reduction, without certifying the subsequent loop-group surjectivity theorem.\par
The split-torus boundary case is consistent with the same hypothesis: \(\mathop{\mathrm{Pic}}(A)=0\) for semilocal \(A\), so every pullback of \(\mathcal O(1)\) by an \(A\)-section is trivial. The map \([x:y]\mapsto[x^d:y^d]\) fixes zero and infinity and pulls \(\mathcal O(1)\) back to \(\mathcal O(d)\). It therefore preserves the endpoint comparison in the reduction, even when \(d\) is not invertible.\par
At the end of the main proof, an open \(C\subset\mathbb A_R^1\) containing an \(R\)-section has finite complement in \(\mathbb P_R^1\) [CF, p. 19]. Each fiber of \(C\) is nonempty, so its closed complement in the corresponding projective line is zero-dimensional. The complement is proper and quasi-finite over \(R\), hence finite. This validates that premise, conditional on the earlier successful reembedding.\par
In Proposition 5.1, the restriction \(H^1(V,Z)\to H^1(R,Z)\) at infinity is surjective: pullback from \(R\) is a right inverse because \(V\) contains the infinity section [CF, p. 20]. The argument from generic vanishing to the clopen neutral-gerbe locus is consistent after decomposition into domains. The cited purity and gerbe equivalence results themselves were not independently proved here.\par
\section*{Conclusion}
One substantive issue is established: the stated rational-point-count condition does not imply the finite-field reembedding assertion. The finite-field replacement using exact residue degrees is verified locally, and two special applications are checked to satisfy it. The property of the external cover in the main proof remains outside the verified scope. No counterexample to the principal Grothendieck–Serre theorem, sectionwise constancy, or simply connected reduction is asserted.\par
Dimension | Recorded state\par
Scientific adjudication | Pending\par
Issue assessment | Major: false reembedding criterion in the core proof chain\par
Repair verification | Finite-field exact-degree criterion and page-11/page-16 uses checked\par
Main-proof repair | Exact-degree property of the cited page-18 cover not independently verified\par
Author adoption | Unknown\par
Independent acceptance | Not performed\par
1 K. Česnavičius and R. Fedorov, Unramified Grothendieck–Serre for isotropic groups, arXiv:2311.08660v3, 21 November 2025, https://arxiv.org/abs/2311.08660v3. Journal identity DOI: https://doi.org/10.4171/jems/1725.\par
\section*{Supporting checks}
The acquired source PDF has 485104 bytes and SHA256\par
a86794f5ee939c60b860ad4639eeebfef069a65eae86f9f6e4e887d2986f17ae.\par
The finite-field counterexample is proved for all finite extensions by [eq:Ycounts] (\#eq:Ycounts)–[eq:Vcounts] (\#eq:Vcounts); no finite numerical search or symbolic program is substituted for that argument.\par
Although the source text was read in full, the cited presentation, isotrivialization, purity, moduli-stack, and loop-group results were not independently consulted. Their appearances in the source do not constitute new verification of them. The absence of a recovered substantive prior report is a bounded recovery result, not a priority claim. No absent computational repository is treated as a mathematical defect.\par
Source acquisition and parsing retain their historical attribution; the specified analytical checks and image views are new. No human review, independent acceptance, or institutional endorsement is claimed.\par
\end{document}
